\documentclass[11pt]{article}

% --- Core Academic Packages ---
\usepackage[utf8]{inputenc}
\usepackage{amsmath, amsthm, amssymb, amsfonts} % Essential for formal econometric math
\usepackage{booktabs}                           % Publication-quality tables (\toprule, \midrule, \bottomrule)
\usepackage{threeparttable}                    % Crucial for robust standard error notes below tables
\usepackage{graphicx}                          % For visual proofs/plots
\usepackage{subcaption}                        % Panel A / Panel B structural layouts
\usepackage{geometry}                           % Standard margins
\usepackage{setspace}                           % Double/1.5 spacing for referee drafts
\usepackage[toc,page]{appendix}                % Formal appendix formatting
\usepackage[colorlinks=true, linkcolor=blue, citecolor=blue]{hyperref}
\usepackage[backend=biber, style=authoryear, maxcitenames=2, uniquename=false]{biblatex}
\usepackage{enumitem}

\geometry{a4paper, margin=0.70in}

\addbibresource{bib.bib}

\title{\textbf{Economic Development, a Political Materialism perspective}}
\author{\textbf{Carlos Gustavo Salas Flores} \\ \small Beijing Normal University}
\date{Oct 2026}

\begin{document}

\maketitle

% =========================================================================
% 1. BACKGROUND
% =========================================================================
\section{Background}
\label{sec:background}

\subsection{The measurement problem in comparative political economy}

Recent economic models attempt to capture latent political variables such as democracy metrics. Some of the most widely used metrics in academia today are V-Dem and Polity-V; however, \textcite{treier_jackman_2008} showed significant error in Polity-V's democracy measures that lead to significant inference errors. In another study, \textcite{little_meng_2024} challenged recent claims of democratic backsliding in recent years and found that such claims are not supported by objective data such as election results, competitiveness of elections, and journalists jailed and murdered; their findings suggest that it is the experts' biases that are changing instead of the objective measures.

These measurement problems matter beyond the democracy literature. If the standard indicators of a country's political character are partly artifacts of expert judgment, then any theory of economic development built on those indicators inherits their errors. This is the entry point of the present work.

\subsection{The institutional theory of development and its limits}

One of the most influential works in recent years has been that of Acemoglu, Johnson, and Robinson (hereafter AJR), who developed a theory of Inclusive and Extractive Institutions that attempts to explain why two territories with similar material conditions can develop very differently, as in the case of Nogales, Arizona, and Nogales, Sonora \parencite{acemoglu_johnson_robinson_2001,acemoglu_johnson_robinson_2005,acemoglu_robinson_2012}. The main drawback of this theory is its inability to develop a theory of the state; instead, it relies on economic indicators and experts' opinions to classify a nation as Inclusive (II) or Extractive (EI). There are, however, at least two flaws in this model. \textit{I) Circular measurement.} Inclusive Institutions are defined by the outcomes they are meant to explain, such as a low expropriation rate and secure property rights. Claiming that a country with Inclusive Institutions displays a low expropriation rate ($A \to X$) licenses the reverse inference that a country displaying a low expropriation rate must have Inclusive Institutions ($X \to A$); the classification then cannot independently explain the outcome, because the outcome was used to construct the classification. \textit{II) Vague horizon.} AJR attempt to answer this objection by arguing that some Extractive Institutions can sustain inclusive economic policies, but that such economies are destined to fail in the long term. Because no time bound is specified, this is an unfalsifiable claim: any continuing success can be reclassified as ``not yet long term.''

A related strand of research links democratization to growth directly. \textcite{acemoglu_naidu_restrepo_robinson_2019} find that transitions to democracy cause subsequent economic growth. Yet this study, too, relies on democracy indicators rather than on a direct analysis of state structures, and the findings of \textcite{little_meng_2024} cast doubt on whether the underlying correlations reflect objective changes in these economies or changes in the biases of the experts who rate them. What is needed, therefore, is a conceptual apparatus that describes domestic politics and international position without depending on contested expert-coded indices.

\subsection{Research question and hypotheses}

This work asks: \emph{which changes in a country's political and international situation actually produce major economic disruption?} To answer it, we must first separate events that the literature habitually lumps together under ``regime change'': the replacement of officeholders, the replacement of the coalition sustaining them, changes in the formal rules of access to power, transformations of the effective structure of political power, restructuring of the state apparatus, policy change, and shifts in a state's external alignment. Section~\ref{sec:framework} develops these distinctions; the hypotheses below anticipate them.

\begin{itemize}
    \item $H_1$ on government turnover: For any two countries $A$ and $B$ with comparable material endowments and similar external alignment, government turnover---and even nominal regime change---that leaves the effective structure of political power untouched will not produce major economic disruption.
    \item $H_2$ on external realignment: For any two such countries, variation in economic disruption tracks variation in external alignment---above all, entry into or exit from a great power's sphere of influence---rather than variation in domestic political events.
    \item $H_3$ on the interaction of domestic and international change: A country that undergoes domestic political change will experience major economic disruption only when that change coincides with an external realignment; and external realignment produces major economic disruption whether or not domestic political change accompanies it.
\end{itemize}

$H_3$ is not a logical consequence of $H_1$ and $H_2$; it is a stronger, separately testable comparative prediction. Section~\ref{sec:hypotheses} states the observable implications of all three hypotheses in a two-by-two design.

% =========================================================================
% 2. CONCEPTUAL FRAMEWORK
% =========================================================================
\section{Conceptual Framework}
\label{sec:framework}

The concepts of the first draft of this work were organized around a triad of \emph{regime}, \emph{establishment}, and \emph{Leviathan}. That triad is abandoned here, because it mixes officeholders, institutions, policies, and the state as a whole. Instead, the framework is organized around five questions: \textbf{who governs}, \textbf{how political power is organized}, \textbf{through which permanent institutions power is exercised}, \textbf{where the state sits internationally}, and \textbf{what outcomes follow}.

\subsection{Units of analysis}

A \emph{country} is the territorial, geographical, demographic, or economic unit; the term is used when comparing geographical conditions or economic performance, not when analyzing the state as a political organization \parencite{linz_stepan_1996,waltz_1979}. The \emph{state} is used in two related senses: domestically, the durable political organization possessing administrative and coercive authority over a territory \parencite{weber1922wirtschaft,hobbes1651leviathan}; internationally, the corporate actor whose alignment and position change relative to other states \parencite{waltz_1979,mearsheimer_2001}. The state is neither the sum of all publicly funded bodies nor the incumbent government. The \emph{international system} is the wider structure within which states operate, characterized by anarchy and by the distribution of capabilities, especially among great powers; the structure constrains behavior but is not itself an actor \parencite{waltz_1979,mearsheimer_2001}. Where cases sit inside a particular regional distribution of power, the relevant unit is the \emph{regional subsystem} \parencite{buzan_waever_2003}.

\subsection{Domestic political elements and actors}

The \emph{government} (or incumbent administration) denotes the persons and organizations currently holding governing office: presidents, prime ministers, ministers, and, where relevant, the ruling coalition. It answers ``who governs?''; policies and international behavior are the government's \emph{actions}, not part of the definition of the actor \parencite{linz_stepan_1996}. The \emph{ruling coalition} is the wider set of parties, elites, and organized groups that sustain the government in power; replacing a president or governing party does not necessarily replace it \parencite{bueno_de_mesquita_2003}.

The \emph{political regime} comprises the rules and institutional arrangements governing access to political power, the distribution of authority among institutions, competition for office, and control or accountability over rulers; it answers ``how is political power formally organized?'' \parencite{linz_stepan_1996,linz_2000}. The \emph{structure of political power}, by contrast, captures the \emph{effective} organization and control of power: the actual separation or concentration of powers, the independence of representation, the party--state relationship, and the capacity of citizens or rival elites to control rulers \parencite{garciatrevijano2010}. A country may change its formal regime while its structure of political power remains intact; Garc{\'i}a-Trevijano's analysis of the party-state (\emph{Estado de partidos}) describes precisely such a configuration, in which state-integrated parties collectively control political power behind formally democratic institutions \parencite{garciatrevijano2010}. \emph{Political parties} are organizations that may participate in government, opposition, or the ruling coalition, and are not automatically part of the government \parencite{sartori_1976}.

The \emph{state apparatus} is the durable administrative, coercive, judicial, and fiscal machinery through which political authority is exercised \parencite{weber1922wirtschaft,althusser_1971}. It replaces the earlier, overly broad notion of an ``establishment'' defined by state funding. Its components are the \emph{coercive apparatus} (armed forces, police, security services; \textcite{tilly_1990,althusser_1971}), the \emph{administrative apparatus} (ministries and public bureaucracy; \textcite{weber1922wirtschaft}), and the \emph{judicial apparatus} (courts and related institutions; \textcite{weber1922wirtschaft,garciatrevijano2010}). Publicly owned or funded bodies such as hospitals, schools, and state companies form a residual \emph{public sector}: membership in it does not automatically make an organization part of the core state apparatus or a strategic political actor.

\subsection{International actors and structural elements}

A \emph{great power} is a state with military capabilities large enough to play a major role in the distribution of power \parencite{mearsheimer_2001}. In asymmetric relationships, the \emph{dominant power} is the stronger state and the \emph{subordinate state} the one whose security, alignment, or foreign-policy choices are substantially constrained by it; international relations frequently exhibit such hierarchy rather than formal equivalence among states \parencite{lake_2009}. A \emph{hegemon} is a state powerful enough to dominate the others in the relevant system or region \parencite{mearsheimer_2001}, while an \emph{aligned state} coordinates its external orientation with another power without that coordination alone proving subordination \parencite{walt_1987}. A \emph{sphere of influence} is a determinate region or relationship in which an external power exercises predominant influence, limiting the independence of the political entities within it, without necessarily implying complete control \parencite{keal_1983,hast_2014}. A state's \emph{external alignment} is its position vis-\`a-vis great powers, alliances, and spheres of influence \parencite{walt_1987}; its \emph{position within the international structure} is its location relative to other powers, so that a change of alignment is a positional change and does not by itself constitute a change in the structure \parencite{waltz_1979}.

\subsection{What actually changes: an event taxonomy}

The literature's undifferentiated use of ``regime change'' is replaced by the following event categories. \emph{Government turnover} replaces the persons or parties holding office \parencite{bueno_de_mesquita_2003}. \emph{Ruling-coalition change} replaces the wider set of actors whose support sustains rule \parencite{bueno_de_mesquita_2003}. \emph{Regime change} alters the formal rules governing access to and organization of political power \parencite{linz_stepan_1996}. A \emph{transformation of the power structure} alters the effective concentration, separation, representation, or control of political power, making it possible to distinguish nominal regime change from deeper structural continuity \parencite{garciatrevijano2010}. \emph{State-apparatus restructuring} transforms the durable administrative, coercive, judicial, or fiscal machinery of the state \parencite{skocpol_1979,althusser_1971}. \emph{Policy change} alters what the government does---its settings, instruments, or overarching goals---and is neither turnover nor regime change \parencite{hall_1993}. Internationally, an \emph{external realignment} changes a state's alliance, dominant-power relationship, or sphere of influence \parencite{walt_1987}, whereas a \emph{systemic power shift} changes the overall distribution of capabilities among states \parencite{gilpin_1981,waltz_1979}. Two distinctions are fundamental: government turnover or even nominal regime change may occur without a transformation of the underlying structure of political power; and a state's external realignment may occur without a systemic power shift, although the two may coincide.

\subsection{Outcomes and conditioning conditions}

The dependent phenomenon is \emph{economic disruption}---which may be differentiated into output collapse, inflation, or capital flight---rather than a static ``economic situation'' \parencite{rodrik_1999}. A state's \emph{security position} and its \emph{military capabilities} are kept analytically separate from economic outcomes \parencite{waltz_1979,mearsheimer_2001}; the first draft's composite ``economic and military situation'' is therefore dropped. Finally, \emph{material endowments} (geography, natural resources, environment, population) are relatively durable conditions that constrain economic and strategic possibilities \parencite{diamond_1997}, while \emph{material capabilities} are the resources a state can actually mobilize---endowments are conditions, capabilities are politically usable capacities \parencite{mearsheimer_2001}.

In compact form: states are governed by incumbent governments, supported by ruling coalitions, operating within a political regime and an effective structure of political power, and acting through a durable state apparatus; internationally, states possess an external alignment and a position within the international structure, especially in relation to great powers and their spheres of influence; and their economic and military outcomes are conditioned by their material endowments and capabilities.

% =========================================================================
% 3. LITERATURE REVIEW AND THEORETICAL FRAMEWORK
% =========================================================================
\section{Literature Review and Theoretical Framework}
\label{sec:lit_review}

Four bodies of literature converge on the research question, and each is incomplete in a way the framework of Section~\ref{sec:framework} addresses.

\paragraph{Institutions and development.} AJR's institutional theory \parencite{acemoglu_johnson_robinson_2001,acemoglu_johnson_robinson_2005,acemoglu_robinson_2012} and the democratization--growth literature \parencite{acemoglu_naidu_restrepo_robinson_2019} explain development through the character of domestic institutions, but measure that character through expert-coded indices whose reliability is contested \parencite{treier_jackman_2008,little_meng_2024}, and---as argued in Section~\ref{sec:background}---through classifications that risk circularity and unfalsifiability. Neither strand distinguishes the replacement of rulers from the transformation of the power structure; both therefore over-attribute economic outcomes to visible domestic political change.

\paragraph{Regimes, power structures, and the state.} The comparative-politics literature supplies the missing domestic distinctions. \textcite{linz_stepan_1996} and \textcite{linz_2000} separate state, regime, and government; \textcite{garciatrevijano2010} adds the effective structure of political power and the party-state paradigm; \textcite{sartori_1976} and \textcite{bueno_de_mesquita_2003} clarify parties and support coalitions; and \textcite{weber1922wirtschaft}, \textcite{althusser_1971}, \textcite{tilly_1990}, and \textcite{skocpol_1979} theorize the durable state apparatus that persists through changes of rulers. \textcite{hall_1993} shows that even policy change comes in distinct orders, so that not every observable policy shift signals deep political change.

\paragraph{International structure and hierarchy.} Structural realism treats states as like units positioned within an anarchic system whose distribution of capabilities constrains them \parencite{waltz_1979}, with great powers competing for power and regional hegemony \parencite{mearsheimer_2001}. Alignment and alliance choices \parencite{walt_1987}, hierarchy between dominant and subordinate states \parencite{lake_2009}, spheres of influence \parencite{keal_1983,hast_2014}, regional security complexes \parencite{buzan_waever_2003}, and systemic power shifts \parencite{gilpin_1981} supply the international half of the framework.

\paragraph{Material conditions and shocks.} Geographical and environmental endowments condition development possibilities \parencite{diamond_1997}, and population and wealth convert into latent and military power \parencite{mearsheimer_2001}. \textcite{rodrik_1999} shows that growth collapses following external shocks are mediated by domestic conflict-management institutions---evidence consistent with the present claim that external position, not domestic political labels, drives disruption.

The gap this thesis occupies is the intersection of these streams: no existing framework explains economic disruption by jointly separating (i) domestic events that change rulers from those that change power structures, and (ii) a state's positional realignment from systemic power shifts. The hypotheses of Section~\ref{sec:background} state exactly this joint claim.

% =========================================================================
% 4. HYPOTHESES AND EMPIRICAL IMPLICATIONS
% =========================================================================
\section{Hypotheses and Empirical Implications}
\label{sec:hypotheses}

The three hypotheses translate into the comparative design of Table~\ref{tab:design}. Cases are selected as pairs of countries with comparable material endowments; the design then varies whether the country experiences (rows) a domestic political change short of power-structure transformation (turnover or nominal regime change) and (columns) an external realignment such as entering or leaving a sphere of influence.

\begin{table}[htbp]
\centering
\begin{threeparttable}
\caption{Predicted economic disruption by type of political and international change}
\label{tab:design}
\begin{tabular}{@{}lcc@{}}
\toprule
 & No external realignment & External realignment \\
\midrule
Government turnover / nominal regime change only & No major disruption ($H_1$) & Major disruption ($H_2$, $H_3$) \\
Power-structure transformation or apparatus restructuring & Ambiguous; out of scope of $H_1$ & Major disruption ($H_2$) \\
No domestic political change & No major disruption & Major disruption ($H_2$, $H_3$) \\
\bottomrule
\end{tabular}
\begin{tablenotes}
\small
\item Note: Predictions assume comparable material endowments across compared cases. ``External realignment'' means a change in alliance, dominant-power relationship, or sphere of influence \parencite{walt_1987,keal_1983}; domestic event categories follow Section~\ref{sec:framework}.
\end{tablenotes}
\end{threeparttable}
\end{table}

Three features make these hypotheses falsifiable where the earlier formulation was not. First, the event categories are defined independently of the outcome, so cases can be classified without circularity. Second, $H_1$ fails if matched cases show major disruption following turnover that left the power structure and external alignment untouched. Third, $H_3$ fails in either direction: if realignment without domestic change produces no disruption, or if domestic change without realignment produces major disruption. The first draft's biconditional (``if and only if it also alters its place in the International Political structure'') asserted a logical consequence of $H_1$ and $H_2$; in fact no such consequence follows from two independence claims, and the reformulated $H_3$ is instead a separate, stronger empirical prediction to be tested on matched cases such as the pairs the first draft gestured at.

% =========================================================================
% 5. CONCLUSIONS & LIMITATIONS
% =========================================================================
\section{Limitations and Conclusion}
\label{sec:discussion}

This draft has rebuilt the conceptual foundations of the thesis. Its main limitations at this stage are four. First, the hypotheses remain untested: the next step is a most-similar-systems case selection among country pairs with comparable material endowments and divergent histories of external realignment. Second, ``major economic disruption'' still requires operationalization (output collapse, inflation, capital flight) along the lines suggested by \textcite{rodrik_1999}. Third, the boundary between nominal regime change and genuine power-structure transformation is analytically clear \parencite{garciatrevijano2010,linz_stepan_1996} but empirically demanding to code; historical cases of state-apparatus restructuring \parencite{skocpol_1979} will guide the coding rules. Fourth, the framework is deliberately agnostic about normative evaluations of regimes: it classifies how power is organized, not whether that organization is desirable.

The core claim survives the revision intact and is now stated precisely: what looks like the economic effect of ``regime change'' is, in many cases, the effect of external realignment---of a state changing its position within the international structure---while domestic political change that leaves the structure of political power untouched leaves the economy largely untouched as well. Whether the data bear this out is the subject of the empirical chapters to come.

\printbibliography

\end{document}
